\section{Core Content}

\subsection{SELECT Essentials}

\subsubsection{The SELECT query mental model (FROM $\rightarrow$ WHERE $\rightarrow$ 
GROUP BY $\rightarrow$ HAVING $\rightarrow$ SELECT $\rightarrow$ ORDER BY $\rightarrow$ LIMIT)}
\J A \texttt{SELECT} query retrieves data from one or more tables. At a minimum, it has:
\begin{itemize}
  \item \texttt{SELECT} (what columns/expressions you want)
  \item \texttt{FROM} (where the data comes from)
\end{itemize}
You can then refine the result using \texttt{WHERE}, \texttt{ORDER BY}, and \texttt{LIMIT/OFFSET}.

\M Even though you write \texttt{SELECT} first, SQL is conceptually processed in this order:
\begin{enumerate}
  \item \texttt{FROM} (choose source tables / build the row set)
  \item \texttt{JOIN ... ON} (combine tables)
  \item \texttt{WHERE} (filter rows)
  \item \texttt{GROUP BY} (group rows)
  \item \texttt{HAVING} (filter groups)
  \item \texttt{SELECT} (compute the output columns)
  \item \texttt{ORDER BY} (sort final result)
  \item \texttt{LIMIT/OFFSET} (return a subset of rows)
\end{enumerate}
This mental model explains many ''weird'' behaviors, like why you can't use a \texttt{SELECT} 
alias inside \texttt{WHERE} in many DBs (because \texttt{WHERE} is evaluated earlier).

\Sx Put filters in \texttt{WHERE} as early as possible to reduce data processed downstream. 
Use \texttt{ORDER BY} only when you need deterministic ordering (otherwise results are not 
guaranteed). Prefer explicit ordering when paginating.

\E The logical order is not the same as the physical execution plan. The optimizer can reorder 
operations, but correctness follows the logical model.

\Pitfall
\begin{itemize}
  \item Assuming results are naturally ordered (they aren't).
  \item Using \texttt{HAVING} for row filters that should be in \texttt{WHERE}.
  \item Confusing \texttt{WHERE} vs \texttt{ON} in joins (changes results for outer joins).
\end{itemize}

\subsubsection{Filtering with WHERE (predicates, boolean logic, precedence)}
\J \texttt{WHERE} filters rows using conditions (predicates), such as comparisons (\texttt{=}, 
\texttt{>}, \texttt{<}), pattern matches (\texttt{LIKE}), and membership checks (\texttt{IN}).

\M SQL boolean logic includes \texttt{AND}, \texttt{OR}, \texttt{NOT}. Operator precedence 
matters: \texttt{NOT} $>$ \texttt{AND} $>$ \texttt{OR}. Use parentheses to make intent explicit. 
For example, ''A and B or C'' is ambiguous; prefer \texttt{(A AND B) OR C} or \texttt{A AND (B OR C)}.

\Sx Make filters sargable (index-friendly): avoid wrapping indexed columns in functions if possible. 
Keep predicates explicit and readable; correctness beats cleverness.

\E PostgreSQL can still use indexes in more cases than you'd expect (expression indexes exist), 
but don't rely on it casually—design queries for clarity first, optimize when needed.

\Pitfall
\begin{itemize}
  \item Missing parentheses causing incorrect results.
  \item Filtering in the wrong place with outer joins (\texttt{WHERE} can turn a 
  \texttt{LEFT JOIN} into an \texttt{INNER JOIN}).
  \item Using \texttt{NOT IN} when NULLs can appear (often wrong results).
\end{itemize}

\subsubsection{NULL semantics (three-valued logic)}
\J \texttt{NULL} means ''unknown / missing''. It is not equal to any value, not even another NULL. 
Use \texttt{IS NULL} and \texttt{IS NOT NULL}.

\M SQL comparisons involving NULL yield \textit{UNKNOWN}, not true/false. Example: 
\texttt{col = NULL} is neither true nor false; it's UNKNOWN, so it won't pass a \texttt{WHERE} 
filter. That's why \texttt{col = NULL} is almost always a bug; you want \texttt{col IS NULL}.

\Sx Prefer \texttt{NOT NULL} constraints where appropriate to reduce NULL complexity. Be careful 
with \texttt{NOT IN (subquery)} if the subquery can return NULL-this can eliminate all rows 
unexpectedly. Prefer \texttt{NOT EXISTS}.

\E NULL handling becomes critical in analytics and outer joins. For reporting, sometimes you use 
\texttt{COALESCE(col, default)} intentionally, but you should understand whether you're replacing 
''unknown'' with a real value.

\Pitfall
\begin{itemize}
  \item Using \texttt{= NULL} instead of \texttt{IS NULL}.
  \item \texttt{NOT IN} with NULLs producing surprising empty results.
  \item Assuming \texttt{COUNT(column)} counts NULLs (it doesn't); use \texttt{COUNT(*)} 
  to count rows.
\end{itemize}

\subsubsection{ORDER BY, LIMIT, OFFSET (determinism and pagination basics)}
\J \texttt{ORDER BY} sorts results. \texttt{LIMIT} restricts the number of rows returned. 
\texttt{OFFSET} skips rows (commonly used for pagination).

\M Without \texttt{ORDER BY}, row order is not guaranteed. Pagination with \texttt{OFFSET} 
can get slower for large offsets because the DB may still scan and discard rows.

\Sx Always define a deterministic ordering for pagination, usually by a unique column 
(\texttt{id}, or \texttt{timestamp + id}). For large datasets, prefer keyset pagination 
(cursor-based) over offset pagination.

\E Keyset pagination typically looks like: \texttt{WHERE (created\_at, id) > 
(:lastCreatedAt, :lastId) ORDER BY created\_at, id LIMIT 50}. This scales better and avoids 
drifting results when data changes.

\newpage
\Pitfall
\begin{itemize}
  \item Paginating without a stable \texttt{ORDER BY} (duplicate/missing rows across pages).
  \item Using OFFSET for very deep pages and expecting good performance.
  \item Sorting on non-indexed columns for large datasets without understanding cost.
\end{itemize}